\subsection{ロピタルの定理の適用条件}
極限を計算していると,分子も分母も $0$ へ近づく商が現れる.このような極限のいくつかは,微分係数の定義に帰着させて求められる.実際,\,$f(c)=g(c)=0$ で,\,$f,g$ が $c$ で微分可能,\,$g'(c)\ne0$ なら,分子と分母を $x-c$ で割ってやれば
\[
 \frac{f(x)}{g(x)}
 =\cfrac{\cfrac{f(x)-f(c)}{x-c}}{\cfrac{g(x)-g(c)}{x-c}}
 \longrightarrow\frac{f'(c)}{g'(c)}
 \qquad(x\to c)
\]
となる.たとえば $f(x)=\e^x-1$,\,$g(x)=\sin x$,\,$c=0$ とすれば,\,$f'(0)=1$,\,$g'(0)=1$ だから $\lim_{x\to0}\frac{\e^x-1}{\sin x}=\frac11=1$ である.

しかし,この方法が使えない場合が3つある.1つ目は,分母が $x^3$ のときのように,\,$g'(c)=0$ となる場合である.2つ目は,\,$x$ を限りなく大きくする極限 $x\to+\infty$ のように,近づく先 $c$ が有限でない場合である.3つ目は,分子と分母がともに $0$ へ近づくのではなく,絶対値が限りなく大きくなる場合,つまり無限大へ発散する場合である.なお,関数についても,数列のときと同じ記号 $+\infty,-\infty$ を使って,\,$x\to+\infty$ や $f(x)\to+\infty$ のように書く.

そこで,1点 $c$ での微分係数の比ではなく,\,$c$ の近くでの導関数の比 $\frac{f'(x)}{g'(x)}$ の極限を調べることを考える.増分の比を導関数の比で表すコーシーの平均値の定理が,ここで役に立つ.ただし,元の商の形を確かめず,機械的に分子と分母を微分してはいけない.次の定理の条件を確かめる必要がある.

\begin{theorem}{ロピタルの定理の片側極限版}\label{定理:ロピタル}
$x$ が区間の一端 $c$ へ近づく片側極限を考える.\,$c$ は有限でも $+\infty,-\infty$ でもよい.端点を除くその区間で $f,g$ は微分可能であり,\,$g'(x)\ne0$ とする.さらに,次のいずれかを満たすとする.
\begin{enumerate}[label=(\roman*),leftmargin=*]
 \item $f(x)\to0$ かつ $g(x)\to0$ である.
 \item $|f(x)|\to\infty$ かつ $|g(x)|\to\infty$ である.
\end{enumerate}
このとき
\[
 \lim_{x\to c}\frac{f'(x)}{g'(x)}=L
\]
が存在すれば
\[
 \lim_{x\to c}\frac{f(x)}{g(x)}=L
\]
である.ここで $L$ は実数のほか,\,$+\infty,-\infty$ でもよい.なお,元の商 $\frac{f(x)}{g(x)}$ は,端点の十分近くで定義されている.
\end{theorem}

(i) を「$0$ と $0$ の形」,(ii) を「無限大と無限大の形」と呼ぶことにしよう.確かめる条件は,商がこのどちらかの形であること,分母の導関数が $0$ にならないこと,導関数の比の極限が存在することの3つである.

両側極限を求めるときは,左右の片側極限にそれぞれ定理を適用し,結果が一致することを確かめる.また,\,$\frac{f'}{g'}$ にもう一度定理を使いたい場合は,その新しい商についても同じ条件を改めて確かめる.微分した後の商が $0$ と $0$ の形になるとは限らないからである.
